\documentclass[11pt]{article}

\usepackage[margin=1in]{geometry}
\usepackage{amsmath,amssymb,amsthm}
\usepackage{graphicx}
\usepackage[hidelinks]{hyperref}
\emergencystretch=3em

\newtheorem{theorem}{Theorem}

\title{A One-Dimensional Obstruction to Unweighted Sampling of Continuous Schauder Frames}
\author{Run fa\_banach\_001, lane 3}
\date{June 21, 2026}

\begin{document}
\maketitle

\section*{Status}

This packet gives a counterexample to Problem~6.1 in Eisner--Freeman,
arXiv:1812.08360 \cite{EF2018}, under the literal unweighted definition of
sampling stated in the source paper. The counterexample is already
one-dimensional.

\section*{Source Problem}

The source paper defines a sampling of a continuous frame
\((x_t,f_t)_{t\in M}\) by choosing a sequence \((t_j)_{j=1}^\infty\) in \(M\)
and considering the discrete sequence \((x_{t_j},f_{t_j})_{j=1}^\infty\).
Problem~6.1 asks which Banach spaces have the property that every bounded
continuous Schauder frame may be sampled to obtain a discrete Schauder frame,
and which Banach spaces admit a bounded continuous Schauder frame which cannot
be so sampled.

\begin{figure}[h]
\centering
\includegraphics[width=\textwidth]{figures/open_problem_crop.png}
\caption{Problem~6.1 from arXiv:1812.08360, final section.}
\end{figure}

\section*{Proof Intuition}

The obstruction is that the source's formal sampling operation records only
the raw sampled pairs. It does not attach quadrature weights. Thus a continuous
identity decomposition over a finite non-atomic measure space can have total
mass one, while every infinite unweighted sample repeats the same rank-one
identity contribution infinitely many times.

This is invisible in weighted Hilbert-frame discretization, where one may
sample with weights. Here, as written, the repeated unweighted coefficient
series diverges.

\begin{theorem}
The one-dimensional Banach space \(\mathbb R\) has a bounded continuous
Schauder frame which cannot be sampled, in the unweighted sense of
\cite{EF2018}, to obtain a discrete Schauder frame.
\end{theorem}

\begin{proof}
Let \(X=\mathbb R\), let \(M=[0,1]\) with Lebesgue measure, and define
\[
  x_t=1\in \mathbb R,\qquad f_t=\operatorname{Id}_{\mathbb R}\in X^*
  \qquad (t\in[0,1]).
\]
The map \(t\mapsto (x_t,f_t)\) is measurable and bounded, with
\(\|x_t\|\,\|f_t\|=1\) for all \(t\).

For every scalar \(x\in\mathbb R\),
\[
  \int_0^1 f_t(x)x_t\,dt
  =
  \int_0^1 x\,dt
  =
  x.
\]
Moreover, for every measurable set \(E\subset[0,1]\), the Pettis integral over
\(E\) is the scalar \(\mu(E)x\). Hence the source definition of continuous
Schauder frame is satisfied.

Now let \((t_j)_{j=1}^\infty\) be any sampling sequence in \(M\). Since the
frame is constant in \(t\), the sampled sequence is
\[
  (x_{t_j},f_{t_j})_{j=1}^\infty
  =
  (1,\operatorname{Id}_{\mathbb R})_{j=1}^\infty .
\]
If this were a discrete Schauder frame for \(\mathbb R\), then for \(x=1\) the
reconstruction series would have to converge to \(1\):
\[
  1
  =
  \sum_{j=1}^{\infty} f_{t_j}(1)x_{t_j}
  =
  \sum_{j=1}^{\infty} 1 .
\]
The series on the right diverges. Therefore no sampling sequence produces a
discrete Schauder frame.
\end{proof}

\section*{Verification Notes}

The example is compatible with the source's continuous Schauder-frame
definition, including the requirement that the Pettis integral over every
measurable subset exist. For each \(E\subset[0,1]\), the partial integral
operator is simply multiplication by \(\mu(E)\).

The example is also bounded in the usual sense relevant to the source's
sampling problem: \(\sup_t\|x_t\|\|f_t\|=1\), and the suppression
unconditionality constant is \(1\).

\section*{Limitations}

This is a counterexample to unweighted sampling exactly as formulated in the
source text. It should not be read as a counterexample to weighted
discretization of continuous Hilbert frames. In the Hilbert-frame literature,
the discretized frame elements are typically allowed to carry sampling weights;
with weights, the present one-dimensional example is harmless.

Thus the mathematical conclusion is best phrased as follows: Problem~6.1 has a
negative answer under the paper's literal unweighted definition, and any
nontrivial positive discretization problem for continuous Schauder frames must
allow weights or impose an additional normalization preventing repeated
unit-mass samples.

\section*{Novelty and Search Bounds}

The local run indexes were searched for arXiv:1812.08360, continuous Schauder
frames, sampling, and discretization. Nearby frame packets concern Schauder
frame false positives and approximate-frame counterexamples, but no packet for
this exact source problem was found.

The source paper itself cites the weighted Hilbert-frame discretization theorem
of Freeman--Speegle. A bounded web search on June~21, 2026 found the source
paper and adjacent continuous-frame discretization material, but no later
Banach-space answer to Problem~6.1 and no note pointing out this unweighted
one-dimensional obstruction.

\section*{Human Review Recommendation}

This packet should be reviewed as a wording-sensitive counterexample. The proof
is elementary; the key review question is interpretive: whether Problem~6.1 is
to be read literally with unweighted raw sampled pairs, or whether the authors
intended weighted sampling by analogy with Hilbert-frame discretization.

\begin{thebibliography}{9}

\bibitem{EF2018}
J.~Eisner and D.~Freeman,
\emph{Continuous Schauder frames for Banach spaces},
arXiv:1812.08360.

\end{thebibliography}

\end{document}
